综合 (3.3.11)--(3.3.14), $\varphi$ 在 $(x_0,H)$ 处的 Jacobian 矩阵具有分块形式
\begin{equation}\tag{3.3.15}
J_{(x_0,H)}(\varphi)=
\left(
\begin{array}{c|c}
0_{1\times n} & 0_{1\times(r-n-1)}\\ \hline
-\operatorname{Hess}(h_{r-n})(0) & 0_{n\times(r-n-1)}\\ \hline
0_{(r-n-1)\times n} & I_{r-n-1}\\ \hline
0_{1\times n} & 0_{1\times(r-n-1)}
\end{array}
\right).
\end{equation}
故 $\varphi$ 在 $(x_0,H)$ 处非分歧, 当且仅当 Hessian 矩阵
\[
\left(
\frac{\partial^2h_{r-n}}{\partial x_i\partial x_j}(0)
\right)_{1\leq i,j\leq n}
\]
可逆. 另一方面, 3.3.2 和 3.3.3 表明 $X\cap H$ 在 $x_0$ 处局部环的完备化为
\begin{equation}\tag{3.3.16}
k[[x_1,\ldots,x_n]]/
\bigl(h_{r-n}(x_1,\ldots,x_n)\bigr),
\end{equation}
结论由定义 1.2 得出.

\subsubsection*{例 3.4.}
假设 $X$ 为 $\mathbb P^r$ 中由齐次方程 $F(X_0,\ldots,X_r)=0$ 定义的超曲面. 则 $\mathbb P(N)=X$, 且在齐次坐标中, $\varphi:X\to\check{\mathbb P}^{r}$ 由下式给出:
\begin{equation}\tag{3.4.1}
\varphi(X_0,\ldots,X_r)=
\left(
\frac{\partial F}{\partial X_0},\ldots,
\frac{\partial F}{\partial X_r}
\right).
\end{equation}
这称为 \underline{Gauss 态射}.

下面是 $\varphi$ 处处分歧的例子. 设 $p$ 为 $k$ 的特征, $r\geq1$, 并取由方程
\begin{equation}\tag{3.4.2}
\sum_{i=0}^{r}X_i^{p^r+1}=0.
\end{equation}
定义的超曲面 $X$. 则
\begin{equation}\tag{3.4.3}
\varphi(X_0,\ldots,X_r)
=(X_0^{p^r},\ldots,X_r^{p^r}).
\end{equation}
所以若 $p=0$, 则 $\check X$ 缩为一个点; 若 $p>0$, 则 $\check X\simeq X$, 且 $\varphi$ 为 Frobenius 态射. 由 3.3(或下面的 3.5.0)可知, 若 $p>2$, 则此嵌入不是 Lefschetz 嵌入.

现在回到一般情形.